\documentclass[../readings.tex]{subfiles}

\begin{document}
\subsection{Numerical Quadrature}
\label{subsec:quadrature}

\textBox{%
Approximating $I = \int_a^b f(x)\,dx$ via weighted sums: $\sum_{i=1}^n w_i f(x_i)$.
}

\subsubsection{Newton-Cotes Rules}

Uses equally-spaced nodes. Higher-order rules are prone to Runge's phenomenon.

\addDef{Trapezoidal Rule}{%
On $n$ panels of width $h = (b-a)/n$:
\begin{align*}
    T_n = h \left[ \frac{f(a) + f(b)}{2} + \sum_{i=1}^{n-1} f(a + ih) \right].
\end{align*}
}
\label{def:trapezoidal-rule}

\addThm{Trapezoidal Error}{%
If $f$ has two continuous derivatives on $[a,b]$, then the composite trapezoidal rule satisfies
\begin{align*}
    I-T_n = O(h^2).
\end{align*}
Thus doubling $n$ reduces the leading error by a factor of about $4$.
}
\label{thm:trapezoidal-error}

\addExcr{%
\begin{enumerate}
    \item Derive the one-panel trapezoidal rule by integrating the linear interpolant through $(a,f(a))$ and $(b,f(b))$.
    \item Use Taylor expansion on one panel to show the local error is $O(h^3)$.
    \item Sum over $O(1/h)$ panels to prove the global $O(h^2)$ statement in Thm~\ref{thm:trapezoidal-error}.
\end{enumerate}
}
\label{ex:trapezoidal-error}

\addDef{Simpson's Rule}{%
Requires $n$ to be even. Uses piecewise quadratic interpolation:
\begin{align*}
    S_n = \frac{h}{3} \left[ f(a) + f(b) + 4\sum_{\text{odd}\;i} f(x_i) + 2\sum_{\text{even}\;i} f(x_i) \right].
\end{align*}
}
\label{def:simpsons-rule}

\addThm{Simpson Exactness and Error}{%
Simpson's rule is exact for polynomials of degree at most $3$. For sufficiently smooth $f$, the composite Simpson rule has global error $O(h^4)$.
}
\label{thm:simpson-error}

\addExcr{%
\begin{enumerate}
    \item Integrate the quadratic interpolant through $x_0,x_1,x_2$ to derive the Simpson weights $1,4,1$.
    \item Verify exactness for $1,x,x^2,x^3$ on a symmetric panel.
    \item Use the exactness degree to explain why the leading error involves the fourth derivative.
    \item Compare the observed error reduction after halving $h$ with Thm~\ref{thm:simpson-error}.
\end{enumerate}
}
\label{ex:simpson-error}

\subsubsection{Gaussian Quadrature}

Optimizes \textbf{both} nodes and weights to achieve the highest possible degree of exactness.

\addThm{Gauss-Legendre Quadrature}{%
Integrating over $[-1, 1]$ with $p$ points:
\begin{align*}
    \int_{-1}^1 f(x)\,dx \approx \sum_{i=1}^p w_i f(x_i).
\end{align*}
The nodes are the roots of the degree-$p$ Legendre polynomial $P_p(x)$, and the rule is exact for all polynomials of degree at most $2p-1$.
}
\label{thm:gauss-legendre}

\addRemark{%
\textbf{(Gauss efficiency)}\enspace Gaussian quadrature spends nodes where orthogonality says they are most useful. For smooth, non-periodic integrands, a small number of Gauss points often beats many equally spaced Newton-Cotes points.
}

\addExcr{%
\begin{enumerate}
    \item For $p=2$, use the roots of $P_2(x)=(3x^2-1)/2$ to find the nodes.
    \item Solve for weights that integrate $1$ and $x^2$ exactly on $[-1,1]$.
    \item Verify that the resulting rule also integrates $x$ and $x^3$ exactly by symmetry.
    \item Explain why exactness through degree $2p-1$ is the maximum possible with $2p$ free parameters.
\end{enumerate}
}
\label{ex:gauss-legendre}

\subsubsection{Richardson Extrapolation and Romberg Integration}

\addThm{Richardson Extrapolation}{%
Given an $O(h^2)$ estimate $A(h)$, build an $O(h^4)$ estimate by:
\begin{align*}
    A_{new} = \frac{4A(h/2) - A(h)}{3}.
\end{align*}
}
\label{thm:richardson}

\addPrf{%
Let $A(h)$ be an $O(h^2)$ approximation of the true value $I$. Using a Taylor expansion in $h$:
\begin{align*}
    A(h) = I + C_1 h^2 + C_2 h^4 + O(h^6).
\end{align*}
Halving the step size gives:
\begin{align*}
    A(h/2) &= I + C_1 (h/2)^2 + C_2 (h/2)^4 + O(h^6) \\
    &= I + \frac{1}{4} C_1 h^2 + \frac{1}{16} C_2 h^4 + O(h^6).
\end{align*}
To eliminate the $O(h^2)$ error term, multiply $A(h/2)$ by 4 and subtract $A(h)$:
\begin{align*}
    4A(h/2) - A(h) &= (4I + C_1 h^2 + \frac{1}{4} C_2 h^4) - (I + C_1 h^2 + C_2 h^4) + O(h^6) \\
    &= 3I - \frac{3}{4} C_2 h^4 + O(h^6).
\end{align*}
Dividing by 3 yields the $O(h^4)$ estimate:
\begin{align*}
    I = \frac{4A(h/2) - A(h)}{3} + O(h^4).
\end{align*}
}

\addDef{Romberg Integration}{%
Iterative application of Richardson extrapolation to the trapezoidal rule.
}
\label{def:romberg}

\addRemark{%
\textbf{(Romberg idea)}\enspace Romberg integration repeatedly cancels the leading even powers in the trapezoidal error expansion. It is effective when $f$ is smooth enough for that expansion to be accurate.
}

\subsubsection{Adaptive Quadrature}

\addRemark{%
\textbf{(Adaptive quadrature)}\enspace Adaptive quadrature focuses function evaluations where $f$ is difficult and uses large panels where $f$ is smooth. The usual mechanism compares one large panel with two smaller panels; if the difference is above tolerance $\varepsilon$, the method recurses.
}

\subsubsection{Periodic Functions and Spectral Accuracy}

\addThm{Periodic Trapezoidal Spectral Accuracy}{%
For smooth periodic functions on $[0,2\pi]$, the trapezoidal rule can converge faster than any algebraic power of $h$. For analytic periodic functions, the convergence is exponential.
}
\label{thm:periodic-trapezoidal}

\addExcr{%
\begin{enumerate}
    \item Apply the trapezoidal rule to $e^{ikx}$ on $[0,2\pi]$ and determine when the discrete sum is exactly zero.
    \item Use a Fourier series argument to explain why smooth periodic functions are integrated unusually well.
    \item Compare errors for $\sin(x)$, $\exp(\cos x)$, and $|x-\pi|$ as $n$ increases.
\end{enumerate}
}
\label{ex:periodic-trapezoidal}

\subsubsection{Exercises}

\addExcr{%
\begin{enumerate}
    \item Doubling nodes in Trapezoidal rule reduces error by $4\times$. Use Thm~\ref{thm:simpson-error} to predict Simpson's reduction factor.
    \item Derive the $p=2$ Gauss points for $[-1, 1]$ using Ex~\ref{ex:gauss-legendre}.
    \item Use Richardson extrapolation from Thm~\ref{thm:richardson} to improve $T(h)$ and $T(h/2)$ to get Simpson's rule.
    \item Integration of $\sin(x)$ on $[0, 2\pi]$ with $n=4$ nodes. Is the error $0$? Why?
\end{enumerate}
}

\end{document}
